\section{The interior mean at a positive cut-off}\label{app:interior}
At a fixed cut-off, \eqref{eq:signed-mean} requires us to sum the logarithms of the row normalizations
up to the stopping time and take the mean of the sum. For each child the first correction has size
\(1/m\). Before summing over \(m\), we must use its cancellation among the children in the same row.
For the four-family predictor, at every regular parent of size \(m\) on the auxiliary macroscopic
event down to \(b\), the insertion formula gives
\[
 L(B+e)-L(B)=\log[p_e(1+\delta_e)]+r_e^{\rm PC},\qquad
 \sum_{e\in i}p_e\delta_e=0,\quad
 |\delta_e|\le C_b/m,\quad |r_e^{\rm PC}|\le C_b/m^2 .
\]
Here the \(p_e\) are the coordinates of the four-family stationary point. Summing over the row gives
\(\log\rho(B,i)=\sum_{e\in i}p_e r_e^{\rm PC}+O_b(m^{-3})\). Replacing \(p_e(1+\delta_e)\) by \(p_e\)
costs \(O_b(m^{-3})\). Write \(V=\ell+\tfrac12\log\det S_{0,\rm gf}+3\sum p_e^2\) for the local
potential of the stationary expansion (\(\ell\) is the local term and the cofactor uses the fixed
gauge). Let \(\delta_e^0\) be its centered linear correction and \(P_e\) the second-order remainder
after extracting \(\Delta_e V=V(B+e)-V(B)\) and \(\delta_e^0\) from \(L(B+e)-L(B)-\log p_e\). Taylor
expansion gives
\[
 r_e^{\rm PC}=\Delta_e V+P_e+\tfrac12(\delta_e^0)^2
                   +(\delta_e^0-\delta_e)+O_b(m^{-3}).
\]
Although \(\delta_e-\delta_e^0=O_b(m^{-2})\) has the order we seek, both corrections have row mean
zero, so their difference vanishes exactly in that mean. The cubic error after summation is
\(O_b(n\sum_{m>j}m^{-3})=O_b(n^{-1})\). The potential contributes its endpoint values. The part we
must average over the whole interval is \(m^2[P_e+(\delta_e^0)^2/2]\). The table pairs each insertion
term with its contributions to the sixteen entries of this physical mean in the continuum.
\begin{center}\small
\begin{tabular}{ll}
\hline
\(-m^2 E_g/2\)& \(G_{\rm cof}\)\\
\(m^2 A_3\)&\(f2E+f2z+f2F\)\\
\(m^2 C^{(3)}\)&\(f2F+C3_{\rm cuad}\)\\
\(-m^2(E_Q+E_R)/2\)&\(f3F+f3c\)\\
\(m^2(-3p_e^2+p_e f_e)\)&\(f4=f4_{\rm expl}+f4_{L0}\)\\
\(m^2\sum_F p_i(1-t_{ie})w_i\)&\(f5c+f5F\)\\
\(-m^2\sum_{E'}p_iq_iw_i\)&\(f6\)\\
\(m^2(\delta_e^0)^2/2\)&\(f7\)\\
first part of \(Q_e\)&\(Q1+\Delta_1\)\\
second part of \(Q_e\), with \(-3m^2\sum_{E'}y_i^2\)&\(Q2+\Delta_2\)\\
\hline
\end{tabular}
\end{center}
\(f2F\) appears twice because the cubic term on the fiber and the pair part of \(C^{(3)}\) have the
same weight. The two parts of \(Q_e\) carry the factors \(-1/2\) and \(1/4\). An available line that
occurs in two factors counts only once, so the correction to the product is
\(\psi\psi'(U_\ell^{-1}-1)\). This gives \(\Delta_1,\Delta_2\). The twelve ordered pairs of distinct
families in the quadratic block contribute \(12/4=3\) times \(m^2\sum_{E'}y_i^2\), canceling the
indicated subtraction at each parent. These identities determine the multiplicities independently of
the numerical evaluation.

The uniform profile error in these scaled means on \([b,1]\) is \(\epsilon_n(b)\to0\). Its total cost
is at most \(\epsilon_n(b)n\sum_{m>j}m^{-2}=o_b(1)\). The predictable indicator of a regular parent is
taken before the row is chosen. Removing the exponentially rare auxiliary exits costs \(o_b(1)\),
since \(n|A|\le C_b\). On \(H_j\), let \(s=A-R_Q(n)+g_T+\Gamma\). There \(|s|=O_b(1/n)\) and
\(\int_{H_j}|d\mu-d\nu|=\mathbb E_\mu[|1-e^{-s-g^\flat}|;H_j]=O_b(1/n)\) by
\eqref{eq:arrival-estimate}. Thus we may use the adaptive law in the Riemann sum. The local potential
telescopes, with scaled endpoint contribution \(\mathcal V(b)/b-\mathcal V(1)\) and error
\(O_b(n^{-1/2})+o_b(1)\).

For the exact rook correction let \(u_e=g_T(M_e)-g_T(M)\), where \(M_e\) here denotes the mask of the
physical child. Thus \(\log\rho^\sharp-\log\rho=\log\mathbb E_{\nu,i}e^{u_e}\). Since
\(|u_e|=O(m^{-2})\), replacing the logarithm by the mean costs \(O(m^{-4})\). The weighted exit costs
and the stopped density give a stopped relative entropy of \(O_b(1)\). The conditional Pinsker
inequality and Cauchy--Schwarz bound the scaled change in the means on passing to \(\mu\) by
\(O_b(n(\sum_{m>j}m^{-4})^{1/2})=O_b(n^{-1/2})\). The increments telescope. The rook endpoints add
\((F(b)-1/12)/b-(F(1)-1/12)\) to the endpoint contribution of the local potential. Together they give
the value of \(G/x\) at \(b\) minus its value at \(1\). Writing this difference as the integral of the
derivative of \(-G/x\) gives the interior mean needed in the final comparison,
\[
 \lim_n n\mathbb E_\mu A^\sharp=\int_b^1 d(x)x^{-2}\,dx+\frac{G(b)}b-G(1)
       =\int_b^1 k(x)x^{-2}\,dx .
\]

\subsection{How to read the row expansion}
The weight of a square depends on how it enters the mean in the preceding table. For terms that
average over a center \(c\), we use the measure \(\nu(c)\,dc\) and write \(\langle\cdot\rangle_c\). A
second independent square \(i\) carries the same measure. If \(i\) lies on a line \(l\) through the
center, only one coordinate is free, either row or column, and the measure on that fiber is
\(d\omega_l(i)=\eta_l\nu(i)\,dy\). In this paragraph \(C\) denotes a kernel distinct from the column
family. Define
\begin{align*}
 A_{lh}&=(L_0-K)_{lh},\qquad B_{lh}=K_{lh}(W_{lh}-K_{lh}),\\
 C_{lh}&=-B_{lh}(W_{lh}-2K_{lh}),\qquad
 w(i)=\sum_{\{l,h\}\subset J(i),\,l\ne h}(L_0)_{lh}.
\end{align*}
Here \(W_{lh}\) equals \(W\) at the crossing, or zero if there is no crossing. We write
\(A(i,c)=\sum_{l\in J(i),\,h\in J(c)}A_{lh}\) and similarly for \(B,C\). If the first argument is a
single line, we sum only over \(h\). After summing, we set \(M_2=A^2+B,\ M_3=A^3+3AB+C\). The cubic
term distinguishes three averages according to whether we choose an independent square, a diagonal
line or a square on a fiber through the center. With their respective signs we obtain
\begin{align*}
 f2E&=\left\langle -M_3(i,c)/6+\tfrac32 W(i)M_2(i,c)\right\rangle_{i,c},\\
 f2z&=-\tfrac16\left\langle
 \int_{l\in D\cup A}\frac{u_l}{U_l}M_3(l,c)\,d\rho_l\right\rangle_c ,\\
 f2F&=\left\langle\sum_{l\in J(c)}\int_{i\in l}W(i)^2\,d\omega_l(i)\right\rangle_c .
\end{align*}
The measure \(d\rho\) is the continuum measure represented by the line weights in
Section~\ref{sec:profile-data}. Four other terms use the same averages.
\begin{align*}
 G_{\rm cof}&=-\tfrac12-3\langle W^2\rangle_c,\qquad
 f4_{\rm expl}=-9\langle W^2\rangle_c,\\
 f4_{L0}&=\langle Ww\rangle_c,\qquad f5c=-3\langle Ww\rangle_c .
\end{align*}
For \(C3_{\rm cuad}\) we consider a center and three translated squares that complete the geometric
square. We take both oriented bases (row--column and diagonal--antidiagonal) and both signs in each
direction. The weight is the product of the four \(W\) factors, \(x^{n_{RC}-3}\) and the \(U_l\)
factors of the distinct diagonal lines. Here \(n_{RC}\) counts the distinct rows and columns. We
integrate with respect to Lebesgue measure \(dr\,dc\,dh\), with \(h>0\) and the square inside the
board. The four \(W\) factors already supply the physical weight, with no additional density at the
center. In the grid implementation a side of length zero carries half weight. In \(f3F\) the quadratic
fiber term includes the variance and the covariance due to shared lines, and \(f3c\) collects the two
center contractions. In \(f5F\) we integrate over the fiber the product of \(-w(i)\) and the sum of
its line responses. In \(f7\) we square \emph{after} summing the responses on the fiber and outside
the center. Finally, \(Q1,Q2\) take the linear and quadratic trace terms with respect to the same
center, with the corrections \(\Delta_1,\Delta_2\) to the product over distinct lines given above.
These measures and multiplicities determine the interior integrand \(d(x)\). At arrival we must
average over the diagonal pairs that remain available in the rook mask.

